% 【本文件写什么】api-v0 契约层，叶子，只写语义不写实现
% - crate 路径：os/components/wateros-mm/mm-frame-alloctor/frame-alloctor-api/api-v0/src/
% - 列出并说明：pub trait、核心类型、错误枚举、常量
% - 每个公开 API 的前置条件、后置条件、线程/中断上下文限制
% - 与 Linux/用户态约定的对齐点与 deliberate 差异
% - v0 版本当前行为与后续可替换 extension 点
% 【事实来源】
% - 上述 crate 下全部 .rs 源文件
% - docs/exports/public-api/ 中对应符号表，以源码为准
% 【不写什么】
% - impl 内部数据结构、汇编、具体算法步骤

\subsubsection{帧分配契约}

\code{PhysicalFrameAllocator} trait 定义帧分配的最小能力：关联类型\code{FrameId}，实现侧与 \code{PhysPageNum} 对齐，以及 \code{alloc\_frame}、\code{dealloc\_frame}。分配粒度为一帧等于一页物理存储，4\,KiB。

\code{FrameAllocError}枚举：\code{OutOfMemory}，池耗尽、\code{InvalidFrame}，释放非法或未分配帧、\code{Unsupported}。\code{mm-api}的\code{MmError::FrameAlloc}经\code{From}映射上述错误。

\code{FrameMemStats}提供\code{total\_frames}、\code{free\_frames}、\code{page\_bytes}只读统计，附带\code{total\_bytes}/\code{free\_bytes}/\code{used\_bytes}便利方法。实现须在关中断或等价临界区内保证单核一致性；本契约不规定 SMP 扩展。

\begin{rustcode}
pub trait PhysicalFrameAllocator {
    type FrameId: Copy + Eq;
    fn alloc_frame(&mut self) -> FrameAllocResult<Self::FrameId>;
    fn dealloc_frame(&mut self, frame: Self::FrameId) -> FrameAllocResult<()>;
}
\end{rustcode}
